\section{Chương 2. Phương pháp đa bước tuyến tính, các loại sai số và tính nhất quán, ổn định, hội tụ}

\subsection{Lưới thời gian, phân loại phương pháp và sai số}
%\begin{itemize}
%    \item Định nghĩa lưới thời gian đều với bước nhảy $h$: $t_n = t_0 + nh$.
%    \item Công thức tổng quát của phương pháp đa bước tuyến tính $k$-bước:
%    \[
%        \sum_{j=0}^{k} \alpha_j y_{n+j} = h \sum_{j=0}^{k} \beta_j f(t_{n+j}, y_{n+j}), \quad \alpha_k \neq 0, \quad |\alpha_0| + |\beta_0| > 0.
%    \]
%\end{itemize}

Để tìm nghiệm xấp xỉ của bài toán giá trị đầu (IVP) trên đoạn $[t_0, T]$, bước đầu tiên là rời rạc hóa miền thời gian liên tục thành một tập hợp hữu hạn các điểm nút.
\subsubsection{Lưới thời gian và công thức tổng quát của phương pháp đa bước}
Ta chia đoạn $[t_0, T]$ thành $N$ đoạn con bằng nhau bởi một \textit{lưới thời gian đều} với bước nhảy $h > 0$:
\begin{equation}
    h = \frac{T - t_0}{N}, \qquad t_n = t_0 + nh, \qquad n = 0, 1, \dots, N. \tag{2.1} \label{eq:2.1}
\end{equation}
\noindent Tại mỗi nút $t_n$, ký hiệu $y(t_n)$ là giá trị nghiệm chính xác (nghiệm giải tích) và $Y_n \approx y(t_n)$ là giá trị nghiệm số xấp xỉ cần tìm. Đồng thời, để viết gọn biểu thức, ta ký hiệu $f_n = f(t_n, y_n)$. 
%Ta bắt đầu từ bài toán (1.1), lấy tích phân hai vế từ $t_n$ đến $t_{n+1}$ ta được:
%\[
%\displaystyle \int_{t_n}^{t_{n+1}} y'(t)dt = \displaystyle \int_{t_n}^{t_{n+1}}f(y;y(t))dt
%\]
%Theo định lý cơ bản của giải tích ra được:
%\[
%y_{n+1}  = y_n +\displaystyle \int_{t_n}^{t_{n+1}}f(y;y(t))dt
%\]
%Tiếp theo ta chuẩn hóa các biến, đặt:
%\[
%t = t_n + sh \qquad 0 \le s\le1 
%\]
%Khi đó:
%\[
%dt = hds
%\]
%Vậy tích phân ở trên trở thành:
%\[
%y_{n+1} = y_n + \displaystyle h\int_{0}^{1}f(t_n + sh))ds
%\]
%Mặt khác, ta dùng nội suy Lagrange cho $f(t_n + sh)$ thì được:
%\[
%f(t_n + sh) \approx P(s) = \displaystyle \sum_{j=0}^{k-1}L_j(s) f_{n-j}
%\]
%Trong đó:
%\[
%L_j(s) = \displaystyle \prod_{m=0\\ m \ne j}^{k-1} \frac{s+m}{m-j}
%\]
%Thay vào tích phân ta được:
%\[
%y_{n+1} = y_n + h \displaystyle h\int_{0}^{1}P(s)ds = y_n + \displaystyle h\int_{0}^{1}\displaystyle %\sum_{j=0}^{k-1}L_j(s) f_{n-j}ds
%\]
%Do tích phân trên là tuyến tính nên:
%\[
%y_{n+1} = y_n + \displaystyle h\sum_{j=0}^{k-1} f_{n-j} \ \displaystyle h\int_{0}^{1}L_j(s)ds
%\]







Một phương pháp đa bước tuyến tính là một công thức sai phân liên hệ trạng thái mới $Y_{n}$ với các trạng thái trong quá khứ thông qua một tổ hợp tuyến tính giữa các giá trị nghiệm và các đạo hàm tương ứng:
\begin{equation}
    \sum_{j=0}^{k} \alpha_j Y_{n+j} = h \sum_{j=0}^{k} \beta_j f(t_{n+j}, Y_{n+j}), \tag{2.2} \label{eq:2.2}
\end{equation}
hoặc viết dưới dạng khai triển tường minh hơn:
\[
    \alpha_k Y_{n+k} + \alpha_{k-1} Y_{n+k-1} + \dots + \alpha_0 Y_n = h \left[ \beta_k f_{n+k} + \beta_{k-1} f_{n+k-1} + \dots + \beta_0 f_n \right].
\]

\begin{baitapbox}{Phương pháp đa bước tuyến tính $k$-bước}
    Một phương pháp đa bước tuyến tính được gọi là $k-$bước khi ít nhất một trong hai hệ số $\alpha_0$ hoặc $\beta_0$ là khác không. Điều này tương đương với:
    \[
    \sum_{j=0}^{k} \alpha_j Y_{n+j} = h \sum_{j=0}^{k} \beta_j f(t_{n+j}, Y_{n+j}) \qquad |\alpha_0| + |\beta_0| > 0
    \]
\end{baitapbox}

\vspace{0.5em}

%\noindent \textbf{Tư duy cấu trúc và điều kiện đặt chỉnh của công thức (2.2):}
%\begin{itemize}
%    \item \textbf{Hệ số chuẩn hóa ($\alpha_k \neq 0$):} Viết phương trình để giải tìm trạng thái mới nhất $y_{n+k}$, do đó hệ số đi kèm với nó bắt buộc phải khác không. Thông thường trong thực hành, ta chuẩn hóa $\alpha_k = 1$.
%    \item \textbf{Độ rộng bước thực sự ($|\alpha_0| + |\beta_0| > 0$):} Điều kiện này bảo đảm rằng mốc thời gian xa nhất trong quá khứ là $t_n$ thực sự tham gia vào công thức. Nếu cả $\alpha_0 = 0$ và $\beta_0 = 0$, công thức tự bị suy biến thành phương pháp $(k-1)$-bước.
%    \item \textbf{Tính tuyến tính:} Công thức được gọi là \textit{tuyến tính} vì nó chỉ là tổ hợp bậc nhất của các $y_{n+j}$ và $f_{n+j}$. Tuy nhiên, bản thân phương trình vi phân $f(t, y)$ hoàn toàn có thể phi tuyến phức tạp theo biến $y$.
%\end{itemize}

\subsubsection{Xây dựng đa thức đặc trưng thứ nhất và thứ hai}

Như đã khảo sát sơ bộ ở phần lý thuyết phương trình sai phân tuyến tính, cấu trúc nghiệm của một dãy truy hồi hệ số hằng được quyết định hoàn toàn bởi các nghiệm của đa thức đại số tương ứng. Đối với phương pháp đa bước tuyến tính tổng quát $k$-bước:
\begin{equation}
    \sum_{j=0}^{k} \alpha_j y_{n+j} = h \sum_{j=0}^{k} \beta_j f(t_{n+j}, y_{n+j}), \tag{2.3} \label{eq:2.3}
\end{equation}
ta có hai tập hợp hệ số độc lập là $\{\alpha_j\}_{j=0}^k$ và $\{\beta_j\}_{j=0}^k$. Để thuận tiện cho việc phân tích độ ổn định và sai số bằng các công cụ giải tích phức, người ta ánh xạ hai bộ hệ số này thành hai đa thức biến phức $\zeta$.

\vspace{0.5em}
\begin{baitapbox}{Đa thức đặc trưng thứ nhất}
    Đa thức này được xây dựng hoàn toàn từ các hệ số vế trái $\alpha_j$ của công thức \eqref{eq:2.3}. Ánh xạ $\rho: \mathbb{C} \to \mathbb{C}$ được định nghĩa là:
\begin{equation}
    \rho(\zeta) = \sum_{j=0}^{k} \alpha_j \zeta^j = \alpha_k \zeta^k + \alpha_{k-1} \zeta^{k-1} + \dots + \alpha_1 \zeta + \alpha_0. \tag{2.4} \label{eq:2.4}
\end{equation}
\end{baitapbox}


%\begin{itemize}
%    \item \textit{Vai trò:} Đa thức $\rho(\zeta)$ chi phối trực tiếp hành vi cơ học của hệ thống khi bước lưới $h \to 0$ (nghĩa là khi vế phải của LMM biến mất, phương trình thoái biến về dạng thuần nhất $\sum \alpha_j y_{n+j} = 0$). 
%    \item Sự phân bố nghiệm của phương trình $\rho(\zeta) = 0$ trên mặt phẳng phức là cơ sở duy nhất để xét \textbf{Điều kiện nghiệm Dahlquist} — yếu tố quyết định tính 0-ổn định của phương pháp. Do giả thiết chuẩn hóa $\alpha_k \neq 0$, đa thức $\rho(\zeta)$ luôn có bậc chính xác là $k$.
%\end{itemize}

\vspace{0.5em}
\begin{baitapbox}{Đa thức đặc trưng thứ hai}
    Đa thức này được xây dựng từ các hệ số vế phải $\beta_j$ của công thức \eqref{eq:2.3}. Ánh xạ $\sigma: \mathbb{C} \to \mathbb{C}$ được định nghĩa là:
\begin{equation}
    \sigma(\zeta) = \sum_{j=0}^{k} \beta_j \zeta^j = \beta_k \zeta^k + \beta_{k-1} \zeta^{k-1} + \dots + \beta_1 \zeta + \beta_0. \tag{2.5} \label{eq:2.5}
\end{equation}
\end{baitapbox}


%\begin{itemize}
%    \item \textit{Vai trò:} Đa thức $\sigma(\zeta)$ đóng vai trò quyết định đến \textbf{độ chính xác (bậc hội tụ)} của phương pháp. 
%    \item Nó kết hợp chặt chẽ với $\rho(\zeta)$ để thiết lập các điều kiện đại số cho tính nhất quán.
%    \item Bậc của $\sigma(\zeta)$ có thể nhỏ hơn $k$ (ví dụ nếu $\beta_k = 0$, ta có phương pháp hiện - Explicit method).
%\end{itemize}

\vspace{0.5em}
%\begin{quote}
%    \textbf{Nhận xét quan trọng:} Do cấu trúc tuyến tính, mọi phương pháp đa bước tuyến tính $k$-bước đều được xác định một cách duy nhất bởi cặp đa thức đặc trưng $(\rho, \sigma)$. Trong văn phong toán học số hiện đại, thay vì liệt kê bảng hệ số, người ta thường gọi tên một phương pháp đa bước đơn giản thông qua bộ đôi đa thức đặc trưng của nó.
%\end{quote}

\subsubsection{Phân loại phương pháp và vấn đề khởi động}
%\begin{itemize}
%    \item \textbf{Phương pháp hiện (Explicit):} $\beta_k = 0$ (Họ Adams-Bashforth).
%    \item \textbf{Phương pháp ẩn (Implicit):} $\beta_k \neq 0$ (Họ Adams-Moulton, BDF) và kỹ thuật toán tử Dự đoán - Hiệu chỉnh (Predictor-Corrector).
%    \item \textbf{Bài toán khởi động lưới:} Giải thích lý do LMM cần $k$ giá trị ban đầu và cách dùng phương pháp một bước (như Runge-Kutta bậc cao) để tạo dữ liệu mồi.
%\end{itemize}

Dựa vào cấu trúc toán học của hệ số $\beta_k$ đi kèm với đạo hàm tại thời điểm đích $t_{n+k}$, phương pháp đa bước tuyến tính được chia làm hai phân lớp mang đặc tính tính toán hoàn toàn khác biệt:
    
    \begin{baitapbox}{Phương pháp hiện}
        Khi $\beta_k = 0$, vế phải của công thức (2.2) chỉ chứa các thông tin ở các thời điểm quá khứ $t_n, t_{n+1}, \dots, t_{n+k-1}$ đã biết hoàn toàn. Giá trị mới $Y_{n+k}$ được tính \textit{tường minh} bằng một công thức đại số trực tiếp:
    \[
        Y_{n+k} = -\sum_{j=0}^{k-1} \frac{\alpha_j}{\alpha_k} Y_{n+j} + h \sum_{j=0}^{k-1}  \frac{\beta_j}{\alpha_k} f_{n+j}.
    \]
    \end{baitapbox}

\textbf{Đại diện:} Họ phương pháp Adams--Bashforth. 

\textbf{Đánh giá:} Tính toán cực kỳ nhanh (chỉ mất $1$ lần tính hàm $f$ cho mỗi bước), nhưng miền ổn định số thường hẹp, dễ mất ổn định khi bước lưới $h$ lớn.
    
    \begin{baitapbox}{Phương pháp ẩn} 
        Khi $\beta_k \neq 0$, ẩn số $Y_{n+k}$ xuất hiện ở \textit{cả hai vế} của phương trình:
    \begin{equation*}
        Y_{n+k} = h \frac{\beta_k}{\alpha_k} f(t_{n+k}, Y_{n+k}) + \text{Các thành phần quá khứ đã biết}.
    \end{equation*}
    \noindent Nếu hàm $f(t, y)$ phi tuyến theo $y$, \eqref{eq:2.3} là một phương trình phi tuyến đối với $Y_{n+k}$. Ta không thể tính trực tiếp mà phải giải phương trình này bằng các kỹ thuật lặp (như lặp điểm bất động hoặc Newton--Raphson). \\
    \end{baitapbox}

\textbf{Đại diện:} Họ phương pháp Adams--Moulton và họ vi phân lùi (BDF). 

\textbf{Đánh giá:} Có vùng ổn định số rộng vượt trội (nhiều công thức BDF đạt tính ổn định vô điều kiện $A$-stability), rất thích hợp cho các hệ phương trình vi phân cứng.\\

\textbf{Kỹ thuật toán tử Dự đoán -- Hiệu chỉnh (Predictor--Corrector)} 

Để tận dụng độ chính xác/tính ổn định cao của phương pháp ẩn mà tránh phải giải phương trình phi tuyến phức tạp với chi phí lớn, trong thực tiễn người ta ghép nối một phương pháp hiện và một phương pháp ẩn thành cặp \textit{Dự đoán -- Hiệu chỉnh} (ký hiệu là sơ đồ PECE):
\begin{itemize}
    \item \textbf{P (Predict):} Dùng một công thức hiện (ví dụ: Adams--Bashforth) để tính giá trị dự đoán sơ bộ $y_{n+k}^{(0)}$.
    \item \textbf{E (Evaluate):} Đánh giá đạo hàm dự đoán $f_{n+k}^{(0)} = f\left(t_{n+k}, y_{n+k}^{(0)}\right)$.
    \item \textbf{C (Correct):} Thay $f_{n+k}^{(0)}$ vào vế phải của công thức ẩn (ví dụ: Adams--Moulton) để sửa lại thành nghiệm hiệu chỉnh chính xác hơn $y_{n+k}^{(1)}$.
    \item \textbf{E (Evaluate):} Đánh giá lại đạo hàm chính thức $f_{n+k}^{(1)} = f\left(t_{n+k}, y_{n+k}^{(1)}\right)$ chuẩn bị cho bước lặp thời gian tiếp theo.
\end{itemize}\\

\noindent\textbf{Bài toán khởi động lưới (The Starting Problem)}
    
Đây là nhược điểm căn bản của mọi phương pháp đa bước. Để tính giá trị đầu tiên do công thức đa bước tạo ra là $Y_k$ (tại $n=0$), thuật toán đòi hỏi phải có đủ $k$ giá trị ban đầu:
\[
    Y_0, \quad Y_1, \quad Y_2, \quad \dots, \quad Y_{k-1}.
\]

Tuy nhiên, bài toán giá trị đầu gốc (IVP) \textbf{chỉ cung cấp duy nhất một điều kiện ban đầu $y_0 = y(t_0)$}. Ta không thể tự ý gán giá trị tùy tiện cho $Y_1, \dots, Y_{k-1}$ vì bất kỳ sai số nào tại bước khởi động sẽ lan truyền và bị khuếch đại làm hỏng toàn bộ quá trình tính toán phía sau.\\

\noindent\textbf{Giải pháp chuẩn tắc trong giải tích số:}

Sử dụng một \textbf{phương pháp một bước có bậc chính xác tương đương hoặc cao hơn} (phổ biến nhất là Runge--Kutta bậc 4 - RK4) xuất phát từ $y_0$ để tính ``mồi'' ra $k-1$ giá trị xấp xỉ chất lượng cao $Y_1, Y_2, \dots, Y_{k-1}$.

Sau khi lưới đã có đủ $k$ nút dữ liệu ban đầu, thuật toán chính thức chuyển sang sử dụng phương pháp đa bước tuyến tính để duy trì tính toán cho tất cả các nút phía sau với tốc độ cao tối đa.


\subsubsection{Toán tử sai phân tuyến tính}
%\begin{itemize}
%    \item Xây dựng toán tử toán học $\mathcal{L}[y](t; h)$ đặc trưng cho cấu trúc sai phân của phương pháp số áp đặt lên nghiệm chính xác.
%\end{itemize}

Để phân tích độ chính xác của phương pháp đa bước tuyến tính $k$-bước một cách độc lập với bài toán vi phân cụ thể, ta liên kết công thức sai phân với một toán tử toán học tác động lên các hàm liên tục khả vi.

\begin{baitapbox}{Toán tử sai phân tuyến tính}
    Với một hàm $y(t)$ bất kỳ khả vi liên tục đủ trơn trên $[t_0, T]$, ta định nghĩa \textit{toán tử sai phân tuyến tính} $\mathcal{L}$ ứng với bước nhảy $h$ như sau:
\begin{equation}
    \mathcal{L}[y](t; h) = \sum_{j=0}^{k} \left[ \alpha_j y(t + jh) - h \beta_j y'(t + jh) \right]. \tag{2.6} \label{eq:2.6}
\end{equation}
    Bằng khai triển Taylor, ta có thể viết lại (2.4) thành một chuỗi lũy thừa theo biến $h$:
\begin{equation}
    \mathcal{L}[y](t; h) = C_0 y(t) + C_1 h y'(t) + C_2 h^2 y''(t) + \dots + C_q h^q y^{(q)}(t) + \dots, \tag{2.7} \label{eq:2.7}
\end{equation}
\end{baitapbox}
% \noindent \textbf{Chứng minh:}

Thật vậy, ta xét khai triển \textit{Taylor} quanh điểm t của $y(t+jh)$ và $y'(t+jh):$
\[
\left\{
\begin{aligned}
    y(t+jh) &= y(t) + jh\cdot y'(t) + \frac{(jh)^2}{2!} \cdot y''(t) + \frac{(jh)^3}{3!} \cdot y'''(t)+...+R_m(t)\\
    y'(t+jh)&= y'(t) + jh \cdot y''(t) + \frac{(jh)^2}{2!} \cdot y'''(t) +...+\frac{m+1}{h} \cdot R_m(t) 
\end{aligned}
\right.
\]

Với $R_m(t) = \dfrac{h^{m+1}}{(m+1)!}\cdot y^{(m+1)}$. Thay vào toán tử thì ta có:
\[
\begin{aligned}
    \mathcal{L}[y](t; h) &= \sum_{j=0}^{k} \left[ \alpha_j y(t + jh) - h \beta_j y'(t + jh) \right]. \\
    \mathcal{L}[y](t; h) &= \sum_{j=0}^{k} \ \alpha_j \left(y(t) + jh\cdot y'(t) + \frac{(jh)^2}{2!} \cdot y''(t) + \frac{(jh)^3}{3!} \cdot y'''(t)+...+R_m(t)\right) \\
    &   \qquad  - h \beta_j \left(y'(t) + jh \cdot y''(t) + \frac{(jh)^2}{2!} \cdot y'''(t) +...+\frac{m+1}{h} \cdot R_m(t)\right) .
\end{aligned}
\]

Từ đây ta thấy toán tử trở thành:
\[
\begin{aligned}
    \mathcal{L}[y](t; h) =  \left(j\alpha_j-\beta_j\right)h\cdot y'(t) &+  \left(\frac{j^2}{2!}\alpha_j - j\beta_j\right)h^2\cdot y''(t) + \left(\frac{j^3}{3!}\alpha_j-\frac{j^2}{2!}\beta_j\right)h^3\cdot y'''(t)+...\\
    &+\left(\frac{j^{m+1}}{(m+1)!}\alpha_j -\frac{j^{m}}{(m)!}\beta_j\right)+...
\end{aligned}
\]

Khi này ta đặt:
\begin{align*}
    C_0 &= \sum_{j=0}^{k} \alpha_j, \\
    C_1 &= \sum_{j=0}^{k} \left( j\alpha_j - \beta_j \right), \\
    C_q &= \sum_{j=0}^{k} \left[ \frac{j^q}{q!} \alpha_j - \frac{j^{q-1}}{(q-1)!} \beta_j \right], \qquad q = 2, 3, \dots
\end{align*}

Thì toán tử trở thành:
\[
    \mathcal{L}[y](t; h) = C_0 y(t) + C_1 h y'(t) + C_2 h^2 y''(t) + \dots + C_q h^q y^{(q)}(t) + \mathcal{O}(h^{p+1}), 
\]

Toán tử $\mathcal{L}[y](t; h)$ đo lường mức độ ``sai lệch'' khi ta áp đặt cấu trúc sai phân của phương pháp số lên một hàm trơn bất kỳ. Phương pháp đa bước được gọi là có \textbf{bậc chính xác $p$} nếu như:
\[
   C_0 = C_1 = C_2 = \dots = C_p = 0 \quad \text{và} \quad C_{p+1} \neq 0.
\]


Khi đó, hằng số $C_{p+1}$ được gọi là \textit{hằng số sai số} (error constant) của phương pháp.

%\noindent trong đó các hệ số hằng $C_q$ hoàn toàn được quyết định bởi các tham số cấu trúc $\{\alpha_j, \beta_j\}$ của phương pháp:

%\begin{align*}
%    C_0 &= \sum_{j=0}^{k} \alpha_j, \\
%    C_1 &= \sum_{j=0}^{k} \left( j\alpha_j - \beta_j \right), \\
%    C_q &= \sum_{j=0}^{k} \left[ \frac{j^q}{q!} \alpha_j - \frac{j^{q-1}}{(q-1)!} \beta_j \right], \qquad q = 2, 3, \dots
%\end{align*}
%\noindent 

\subsubsection{Phân loại và phân tích các nguồn sai số trong tính toán số}
%\begin{itemize}
%    \item \textbf{Sai số chặt cụt cục bộ (Local Truncation Error - LTE):} Định nghĩa toán học thông qua toán tử sai phân khi giả định nghiệm quá khứ hoàn toàn chính xác; kỹ thuật xác định hằng số sai số mạng.
%    \item \textbf{Sai số toàn cục (Global Error):} Sự tích lũy và lan truyền của các sai số cục bộ qua từng bước tính trên toàn bộ khoảng thời gian $e_n = y(t_n) - y_n$.
%    \item \textbf{Sai số làm tròn máy tính (Round-off Error):} Sai số do giới hạn biểu diễn số thực trên máy tính (băm nhỏ số); tác động triệt tiêu ưu thế của việc chọn bước lưới $h$ quá nhỏ.
%    \item \textbf{Sai số dữ liệu ban đầu / Sai số nhiễu (Inherited Error):} Sai số lan truyền xuất phát từ sự sai lệch nhỏ trong việc chuẩn bị các giá trị khởi động $y_1, y_2, \dots, y_{k-1}$.
%\end{itemize}

Trong quá trình triển khai phương pháp đa bước trên máy tính thực tế, nghiệm số $Y_n$ thu được luôn bị sai lệch so với nghiệm giải tích $y(t_n)$. Tổng sai số này là sự kết hợp và cạnh tranh của $4$ nguồn sai số cơ bản dưới đây.
\begin{baitapbox}{Sai số cục bộ}
    Sai số cục bộ là lượng sai số phát sinh sinh ra \textit{chỉ sau một bước tính toán duy nhất}, dưới giả định lý thuyết rằng tất cả $k$ giá trị quá khứ được đưa vào công thức đều hoàn toàn chính xác, sai số cục bộ được tính bởi:
    \[  
        E_n= y(t_n) - Y_n \qquad Y_s = y(t_s) \  \forall s =\overline{0,...,k}
    \]
\end{baitapbox}
Bằng định nghĩa trên của sai số chặt cụt cục bộ tại bước thứ $n+k$ là toán tử sai phân tác động lên nghiệm giải tích đúng $y(t)$, được chuẩn hóa bởi hệ số $\alpha_k$:
\begin{baitapbox}{Sai số cục bộ của phương pháp đa bước}
    Cho bài toán (IVP) thỏa $f$ khả vi liên tục và gọi $y(t)$ là nghiệm chính xác của nó.Đại lượng \textit{sai số cục bộ} được cho bởi:
    \begin{equation}
        E_{n} = \frac{1}{\alpha_n} \mathcal{L}[y](t_n; h). \tag{2.8} \label{eq:2.8}
    \end{equation}   
\end{baitapbox}
%\noindent\textbf{Chứng minh:}

Ta xét biểu thức tính sai số cục bộ:
\[
E_n = y(t_n) - Y_n
\]

Ta giả sử nghiệm số $Y_n$ được xác định bởi phương pháp đa bước dạng ẩn bởi phương trình:
\[
\sum_{j=0}^{n} \left[ \alpha_j Y_{j}  -h  \beta_j f(t_{j}, Y_{j})\right] = 0
\]

Ta giả sử toàn bộ các giá trị quá khứ $Y_{n-i}$ ($\ \forall i = \overline{1,n-1}$) chính xác tuyệt đối, tức là $Y_i = y(t_i) \ \forall i  = \overline{1,n-1}$, khi này biểu thức tính $Y_n$ trở thành:
\[
\Leftrightarrow \sum_{j=0}^{n-1} \left[ \alpha_j y(t_j)  -h  \beta_j f(t_{j}, y(t_j))\right] + \alpha_n Y_{n}  -h  \beta_n f(t_{n}, Y_{n})  = 0
\]
\[
\Leftrightarrow  \sum_{j=0}^{n} \left[ \alpha_j y(t_j)  -h  \beta_j f(t_{j}, y(t_j))\right] + \alpha_n Y_{n}  -h  \beta_n f(t_{n}, Y_{n}) -\left[ \alpha_n y(t_n)  -h  \beta_n f(t_{n}, y(t_n)) \right] = 0
\]

Ta nhận thấy cụm sigma chính là toán tử sai phân tuyến tính, thay vào phương trình ta được:
\[
\mathcal{L}[y](t; h) = \alpha_n \left[ y(t_n) - Y_n\right] + h\beta_n\left[f(t_n;y(t_n)) -f(t_n;Y_n) \right]
\]

Áp dụng định lý giá trị trung bình cho hàm $f$ theo biến $y$, ta được:
\[
f(t_n;y(t_n)) -f(t_n;Y_n) = \frac{\partial f}{\partial y}(t_n;\eta)\left[y(t_n) - Y_n\right] \quad (\text{Với} \ \eta \ \text{nằm giữa} \  y(t_n) \text{và} \ Y_n)
\]

Thay vào phương trình, ta có:
\[
\mathcal{L}[y](t; h) = \alpha_n \left[ y(t_n) - Y_n\right] - h\beta_n \frac{\partial f}{\partial y}(t_n;\eta)\left[y(t_n) - Y_n\right]
\]

Ta xét bước lưới tiến về $0$, tức $h \to 0$, thì được:
\[
\mathcal{L}[y](t; h) = \alpha_n \left[ y(t_n) - Y_n\right]
\]
\[
\Leftrightarrow E_n = \frac{1}{\alpha_n} \mathcal{L}[y](t_n; h)
\]

%\noindent Nếu phương pháp đạt bậc chính xác $p$, áp dụng khai triển (2.5) cho nghiệm đúng $y(t)$ (thỏa mãn $y'(t) = f(t, y)$), biểu thức chủ đạo của sai số chặt cụt cục bộ khi $h \to 0$ là:
%\begin{equation}
%    E_{n+k} = \frac{C_{p+1}}{\alpha_k} h^{p+1} y^{(p+1)}(t_n) + \mathcal{O}(h^{p+2}). \tag{2.7}
%\end{equation}
%\noindent \textbf{Ý nghĩa:} Tốc độ tiến về $0$ của LTE tỷ lệ thuận với $h^{p+1}$. Hằng số $\frac{C_{p+1}}{\alpha_k}$ là chỉ số quan trọng để so sánh độ chính xác giữa các phương pháp có cùng bậc $p$.
\begin{baitapbox}{Sai số chặt cụt cục bộ của phương pháp}
    Sai số chặt cụt của phương pháp đa bước được cho bởi:
    \begin{equation}
         T_n = \frac{ \sum_{j=0}^{k} \left[ \alpha_j y(t_{n+j})  -h  \beta_j f(t_{n+j}, y(t_{n+j}))\right]}{h \sum_{j=0}^{k} \beta_j} = \frac{\mathcal{L}[y](t_n; h)}{h \sigma(1)} \tag{2.9} \label{eq:2.9}
    \end{equation}
   

\end{baitapbox}
%\noindent \textbf{Chứng minh:}


Thật vậy, ta xét bài toán (IVP) và phương pháp đa bước có dạng:
\[
y'=f(t;y) \qquad \sum_{j=0}^{n} \alpha_j Y_{n+j}  =h  \sum_{j=0}^{n}\beta_j f(t_{n+j}, Y_j)
\]

Sai số cắt cục đo lường mức độ sai lệch của công thức xấp xỉ khi ta được \textit{nghiệm chính xác} vào công thức của phương pháp số. Ta giả sử (IVP) có nghiệm chính xác $y(t)$ thỏa $y'=f(t;y)$. Thay vào công thức tổng quát phương pháp đa bước ta tính được phần dư:
\[
R_n = \sum_{j=0}^{k} \alpha_j y(t_{n+j})  -h  \sum_{j=0}^{k}\beta_j f(t_{j}, y(t_{n+j}))
\]

Bằng công thứ sai phân tuyến tính ta có phần dư có thể được viết lại:
\[
R_n = \mathcal{L}[y](t_n; h) = \sum \alpha_j y(t) + \left(\sum j\alpha_j -\sum \beta_j\right) h y'(t) + \mathcal{O}({h^2})
\]

Ta thu phóng phần dư này một cách thích hợp để $T_n$ tương tự như $y'-f(t;y)$, trong trường hợp này ta chuẩn hóa bằng chia các vế cho $h\sum\beta_j$:
\[
T_n = \frac{R_n}{h\sum \beta_j} = \frac{\mathcal{L}[y](t_n; h)}{h\sigma(1)}
\]
% \textbf{Sai số toàn cục (Global Error)}\\
%    Trong tính toán lặp, ta không bao giờ có dữ liệu quá khứ chính xác tuyệt đối. Sai số chặt cụt cục bộ $d_{n+k}$ tại mỗi bước sẽ không mất đi mà bị \textit{tích lũy và lan truyền} qua các bước tính tiếp theo trên toàn khoảng $[t_0, T]$. 

%Sai số toàn cục tại nút $t_n$ được định nghĩa là độ lệch thực tế giữa nghiệm đúng và nghiệm số:
%\begin{equation}
%    e_n = Y(t_n) - y(t_n). \tag{2.8}
%\end{equation}
%\noindent Mặc dù sai số cục bộ ở mỗi bước có bậc $\mathcal{O}(h^{p+1})$, do số bước lặp cần thực hiện trên đoạn $[t_0, T]$ là $N = \frac{T - t_0}{h} \sim \mathcal{O}(h^{-1})$, tổng tích lũy của chúng khiến sai số toàn cục bị giảm đi một bậc, chỉ còn đạt %$\mathcal{O}(h^p)$:
%\[
%    \|e_n\| \le K h^p \qquad (K > 0 \text{ là hằng số độc lập với } h).
%\]

%    \item \textbf{Sai số làm tròn máy tính (Round-off Error)}\\
%    Sai số làm tròn phát sinh do hệ thống máy tính chỉ có khả năng biểu diễn số thực trên một số lượng bit hữu hạn (ví dụ: chuẩn IEEE 754 số dấu phẩy động 64-bit đôi chỉ giữ lại khoảng 15--17 chữ số thập phân có nghĩa). 

%Gọi $Y_n$ là nghiệm thực tế máy tính tính ra (đã bao hàm làm tròn). Tại mỗi bước tính, phép tính số học tạo ra một nhiễu làm tròn $\rho_n$, sao cho:
%\[
%    \sum_{j=0}^{k} \alpha_j Y_{n+j} = h \sum_{j=0}^{k} \beta_j f(t_{n+j}, Y_{n+j}) + \rho_{n+k}.
%\]
%\noindent \textbf{Nghịch lý cạnh tranh sai số:} Để giảm sai số chặt cụt toàn cục ($\sim h^p$), ta muốn chọn bước lưới $h$ cực kỳ nhỏ. Tuy nhiên, khi $h \to 0$, số bước lặp $N = \frac{T-t_0}{h} \to \infty$, khiến tổng sai số làm tròn tích lũy ($\sim \frac{\epsilon_{\text{machine}}}{h}$) bùng nổ vô cùng nhanh. Thực tế luôn tồn tại một bước lưới tối ưu $h_{\text{opt}}$ mà tại đó tổng sai số là nhỏ nhất; việc băm nhỏ $h < h_{\text{opt}}$ sẽ phản tác dụng và làm độ chính xác xấu đi rệt.

% \textbf{Sai số dữ liệu ban đầu / Sai số nhiễu (Inherited Error)}\\
%     Phương pháp đa bước $k$-bước đòi hỏi $k$ giá trị xuất phát $y_0, y_1, \dots, y_{k-1}$. Trong thực tế:
% \begin{itemize}
%     \item Bản thân điều kiện đầu gốc $y_0$ có thể đã mang sai số đo đạc thực nghiệm.
%     \item Các giá trị mồi $y_1, \dots, y_{k-1}$ được tạo ra từ một phương pháp một bước (như RK4) cũng chứa sẵn sai số xấp xỉ của phương pháp đó.
% \end{itemize}
% \noindent Sai số ban đầu $\delta_j = y(t_j) - y_j$ ($j = 0, \dots, k-1$) được gọi là \textit{sai số kế thừa}. Hành vi lan truyền của tập nhiễu ban đầu này qua phương trình sai phân đa bước chính là đối tượng nghiên cứu cốt lõi của lý thuyết \textbf{Tính ổn định số}. Nếu phương pháp không ổn định, một sai số $\delta_j$ dù nhỏ tùy ý cũng sẽ bị khuếch đại hàm mũ theo thời gian, phá hủy hoàn toàn nghiệm số.



% ==========================================
\subsection{Tính nhất quán, tính ổn định và sự hội tụ của phương pháp đa bước tuyến tính}
 
\subsubsection{Tính nhất quán}
%\begin{itemize}
%    \item \textbf{Định nghĩa:} Phương pháp số được gọi là nhất quán nếu sai số chặt cụt cục bộ tiến về $0$ khi bước lưới $h \to 0$.
%    \item \textbf{Điều kiện đại số:} Thiết lập thông qua đa thức đặc trưng: $\rho(1) = 0$ và $\rho'(1) = \sigma(1)$. Điều này tương đương phương pháp đạt bậc chính xác tối thiểu $p \ge 1$.
%\end{itemize}

Để một phương pháp xấp xỉ số có ý nghĩa thực tiễn, yêu cầu tối thiểu đầu tiên là cấu trúc sai phân của nó phải ``mô phỏng đúng'' phương trình vi phân ban đầu khi ta thu nhỏ kích thước bước lưới. Tính chất cơ bản này được gọi là \textit{tính nhất quán}.

\vspace{0.5em}
\begin{baitapbox}{\textbf{Tính nhất quán}}
    Một phương pháp đa bước tuyến tính $k$-bước được gọi là \textbf{nhất quán} với phương trình vi phân $y'(t) = f(t, y(t))$ nếu sai số chặt cụt $T_{n}$ tiến về $0$ một cách đồng đều khi bước nhảy lưới $h \to 0$, tức là:
    \begin{equation}
        \lim_{h \to 0} \max_{0 \le n \le N} \|T_{n}\| = 0, \tag{2.10} \label{eq:2.10}
    \end{equation}
    \noindent với giả định hàm nghiệm $y(t)$ là hàm khả vi liên tục đủ trơn trên miền khảo sát.
\end{baitapbox}\\

\vspace{0.5em}
\begin{baitapbox}{\textbf{Điều kiện đại số thông qua đa thức đặc trưng}}
    \noindent Để tiện cho việc phân tích đại số, ta nhắc lại cặp \textit{đa thức đặc trưng thứ nhất} $\rho(\zeta)$ và \textit{đa thức đặc trưng thứ hai} $\sigma(\zeta)$ gắn liền với các hệ số $\{\alpha_j, \beta_j\}$ của phương pháp đa bước:
\begin{equation}
    \rho(\zeta) = \sum_{j=0}^{k} \alpha_j \zeta^j, \qquad \sigma(\zeta) = \sum_{j=0}^{k} \beta_j \zeta^j. \tag{2.11} \label{eq:2.11}
\end{equation}
Một phương pháp đa bước tuyến tính nhất quán khi và chỉ khi đồng thời thỏa mãn hai điều kiện đại số:
\begin{equation}
    \rho(1) = 0 \qquad \text{và} \qquad \rho'(1) = \sigma(1). \tag{2.12} \label{eq:2.12}
\end{equation}
\noindent Điều kiện này tương đương với việc phương pháp đạt bậc chính xác tối thiểu $p \ge 1$. Khi tính nhất quán được thỏa mãn, sai số chặt cụt cục bộ sẽ tiến về $0$ với tốc độ ít nhất là tuyến tính: $T_{n} = \mathcal{O}(h)$.
\end{baitapbox} 

\noindent\textbf{Chứng minh:}

Từ công thức sai số chặt cụt \eqref{eq:2.7} và toán tử sai phân tuyến tính \eqref{eq:2.9}, ta có phương trình:
\[
    T_{n} = \frac{\mathcal{L}[y](t_n; h)}{h \sigma(1)} = \frac{1}{h\sigma(1)}\left[C_0 y(t) + C_1 h y'(t) + C_2 h^2 y''(t) + \dots + C_q h^q y^{(q)}(t) +...\right]
\]


Do đó, điều kiện \eqref{eq:2.10} đòi hỏi khi $h \to 0$, số hạng tự do không chứa $h$ và số hạng chứa bậc $h^{-1}$ phải triệt tiêu, nghĩa là phương pháp bắt buộc phải thỏa mãn:
\[
    C_0 = 0 \quad \text{và} \quad C_1 = 0. 
\]
\noindent\textbf{Biểu diễn $C_0$ qua đa thức đặc trưng:} Theo định nghĩa của hệ số Taylor, ta có:
    \[
        C_0 = \sum_{j=0}^{k} \alpha_j = \rho(1).
    \]

Do đó, điều kiện $C_0 = 0$ hoàn toàn tương đương với phương trình đại số:
\begin{equation}
      \rho(1) = 0. \tag{2.13} \label{eq:2.13}
\end{equation}

\textit{Ý nghĩa:} Phương trình $\rho(1) = 0$ bảo đảm rằng nếu bài toán có nghiệm là hàm hằng ($y(t) \equiv \text{const} \Rightarrow y' = 0$), thì phương pháp đa bước tính toán ra kết quả chính xác tuyệt đối. Nói cách khác, $\zeta = 1$ luôn phải là một nghiệm của đa thức đặc trưng thứ nhất $\rho(\zeta)$ (nghiệm này được gọi là \textit{nghiệm chính}).\\

\noindent \textbf{Biểu diễn $C_1$ qua đa thức đặc trưng:} Đạo hàm bậc nhất của $\rho(\zeta)$ tại điểm $\zeta = 1$ là:
\[
    \rho'(\zeta) = \sum_{j=1}^{k} j \alpha_j \zeta^{j-1} \implies \rho'(1) = \sum_{j=0}^{k} j \alpha_j.
\]

Mặt khác, giá trị của đa thức đặc trưng thứ hai tại $\zeta = 1$ là:
\[
    \sigma(1) = \sum_{j=0}^{k} \beta_j.
\]

Theo định nghĩa của $C_1$, ta có:
\[
    C_1 = \sum_{j=0}^{k} j \alpha_j - \sum_{j=0}^{k} \beta_j = \rho'(1) - \sigma(1).
\]

Vì vậy, điều kiện $C_1 = 0$ hoàn toàn tương đương với phương trình đại số:
\begin{equation}
    \rho'(1) = \sigma(1). \tag{2.14} \label{eq:2.14} 
\end{equation}

\textit{Ý nghĩa:} Phương trình $\rho'(1) = \sigma(1)$ bảo đảm rằng phương pháp xấp xỉ chính xác cho các nghiệm dạng tuyến tính ($y(t) = ct$).




\subsubsection{Tính ổn định tổng quát trong lan truyền nhiễu}
%\begin{itemize}
%    \item Định nghĩa tính ổn định tổng quát: Khả năng một phương pháp số giới hạn biên độ khuếch đại của các loại sai số (sai số làm tròn, sai số dữ liệu ban đầu) không bùng nổ vô hạn khi thực hiện lặp số.
%\end{itemize}

Một phương pháp số có thể nhất quán (đạt độ chính xác cao trên từng bước độc lập) nhưng vẫn thất bại hoàn toàn trong thực tiễn tính toán nếu nó không có khả năng kiểm soát sự lan truyền của các nhiễu nhỏ. Khả năng kiềm chế sự khuếch đại sai số này được định nghĩa là \textit{tính ổn định}.

\vspace{0.5em}
%\noindent \textbf{Tính ổn định của phương pháp số}
%\begin{itemize}
%    \item \textbf{Khái niệm cốt lõi:} Khả năng một phương pháp số giới hạn biên độ khuếch đại của các loại sai số (như sai số làm tròn máy tính, sai số dữ liệu ban đầu hay sai số chặt cụt tại mỗi bước lặp) không bị bùng nổ vô hạn khi thực hiện quá trình lặp số trên một miền thời gian $[t_0, T]$ cố định.
%\end{itemize}
\begin{baitapbox}{Tính 0-ổn định}
    Phương pháp đa bước tuyến tính $k$-bước được gọi là \textbf{0-ổn định} nếu tồn tại một hằng số $K > 0$ (độc lập với bước lưới $h$) sao cho đối với hai dãy nghiệm bất kỳ $\{Y_n\}$ và $\{Z_n\}$ được sinh ra từ cùng một phương pháp nhưng với dữ liệu ban đầu khác nhau:
    \begin{align*}
        \sum_{j=0}^{k} \alpha_j Y_{n+j} &= h \sum_{j=0}^{k} \beta_j f(t_{n+j}, y_{n+j})  \\
        \sum_{j=0}^{k} \alpha_j Z_{n+j} &= h \sum_{j=0}^{k} \beta_j f(t_{n+j}, z_{n+j}) 
    \end{align*}
    \noindent ta luôn có sự bị chặn của độ lệch toàn cục:
    \begin{equation}
        |Y_n - Z_n| \le K  \max_{0 \le j \le k-1} |Y_j - Z_j| , \tag{2.15} \label{eq:2.15}
    \end{equation}
    \noindent với mọi $n = k, k+1, \dots, N$ sao cho $t_n \in [t_0, T]$
\end{baitapbox}
\vspace{0.5em}
\noindent \textbf{Điều kiện nghiệm Dahlquist}

Xét công thức tổng quát của phương pháp đa bước, khi bước lưới tiến dần đến không ($h \to 0$).
\[
 \sum_{j=0}^{k} \alpha_j Y_{n+j} = h \sum_{j=0}^{k} \beta_j f(t_{n+j}, y_{n+j}) \to 0
\]

Khi đó, phương trình vi phân giảm cấp về dạng tầm thường $y' = 0$, và công thức đa bước suy biến thành phương trình sai phân thuần nhất:
\[
    \alpha_k y_{n+k} + \alpha_{k-1} y_{n+k-1} + \dots + \alpha_0 y_n = 0.
\]

Để nghiệm của phương trình sai phân này không bị bùng nổ hàm mũ khi $n \to \infty$, nhà toán học Germund Dahlquist đã chứng minh điều kiện cần và đủ cho tính 0-ổn định:

\begin{baitapbox}{Định lý Dahlquist về tính 0-ổn định:} 
    Một phương pháp đa bước tuyến tính là \textbf{0-ổn định} khi và chỉ khi mọi nghiệm $\zeta_i$ của đa thức đặc trưng thứ nhất $\rho(\zeta) = \sum_{j=0}^k \alpha_j \zeta^j = 0$ đều thỏa mãn \textit{Điều kiện nghiệm}:
    \begin{enumerate}
        \item Tất cả các nghiệm đều nằm trong hoặc trên đường tròn đơn vị: $|\zeta_i| \le 1, \;\forall i$.
        \item Nếu có nghiệm nằm trực tiếp trên đường tròn đơn vị ($|\zeta_i| = 1$), thì nghiệm đó buộc phải là \textbf{nghiệm đơn} (bội bội bằng $1$).
    \end{enumerate}
\end{baitapbox}
\noindent \textbf{Chứng minh:}
\noindent\textbf{Chiều thuận:}$0-$ổn định $\Rightarrow$ điều kiện nghiệm
Ta xét phương pháp đa bước tuyến tính $k$-bước tổng quát:
\[
\sum_{j=0}^{k} \alpha_j Y_{n+j} = h \sum_{j=0}^{k} \beta_j f(t_{n+j}, y_{n+j})
\]

Ta áp dụng bài toán $y'=f(t;y(t) = 0$ cho phương pháp trên thì được:
\[
\sum_{j=0}^{k} \alpha_j Y_{n+j} = 0
\]

Đây là phương trình sai phân tuyến tính hệ số hằng có nghiệm tổng quát là:
\[
Y_n = \sum_{s} p_s(n) \zeta_s^n
\]

Trong đó: 
\begin{itemize}[label = \null]
    \item $\zeta_s$ là nghiệm của đa thức đặc trưng thứ nhất $\rho(\zeta)$.
    \item $p_s(n)$ là đa thức theo $n$ biến, bậc của đa thức nhỏ hơn bậc của $\zeta_s$ 1 đơn vị.
\end{itemize}
Ta giả sử nghiệm của bài toán trên không thỏa mãn điều kiện nghiệm, tức là:
\begin{enumerate}
    \item Tồn tại nghiệm nằm ngoài đường tròn đơn vị : $ \exists \ i$ sao cho $|\zeta_i| > 1  $
    \item Nếu có nghiệm nằm trên đường tròn đơn vị ($|\zeta_i| = 1  $) thì đó là nghiệm bội $n$ \quad $n \ge 2$.
\end{enumerate}

Đối với các nghiệm nằm ngoài đường tròn tức $|\zeta_i| > 1$, khi này:
\[Y_n = \sum_{s} p_s(n) \zeta_s^n \to \infty \quad  \text{khi} \quad n \to \infty\]

Vậy ta có độ lệch toàn cục không bị chặn $\Rightarrow$ bài toán không có tính $0$-ổn định.

Đối với các nghiệm bội $m_s$ thì đa thức $p_s(n)$ có bậc là $m_s -1 \ge 1$, khi này đa thức $p_s(n)$ có tốc độ tăng trưởng là $n^{m_s -1}$, khi $n \to \infty$ thì $p_s(n) \to \infty$ tức là:
\[
Y_n = \sum_{s} p_s(n) \zeta_s^n \to \infty 
\]

Vậy tương tự như trên ta có bài toán không có tính $0$-ổn đinh. Theo mệnh đề phản đảo, ta có được một bài toán có tính $0$-ổn định thì nó phải thỏa mãn điều kiện nghiệm.\\

\noindent\textbf{Chiều đảo:} $0-$ổn định $\Leftarrow$ điều kiện nghiệm.
Đối với chiều thuận chứng minh khá dài và nặng về mặt kỹ thuật nên phần này có thể bỏ qua trong báo cáo này. Qua đó có thể tham khảo phác thảo chứng minh sau:
Chiều đảo được chứng minh bằng cách chỉ ra rằng khi thỏa mãn điều kiện nghiệm, tất cả các nghiệm của hệ thức truy hồi đều bị chặn. Nghiệm tổng quát là tổ hợp tuyến tính của các dãy cơ sở gắn liền với từng nghiệm $z_\alpha$ có bội $m_\alpha$: với $k = 0, \dots, m_\alpha - 1$, các số hạng có dạng
\[
    y_n = \left. \frac{d^k}{dz^k} (z^n) \right|_{z=z_\alpha},
\]
dạng này khái quát hóa sự tăng trưởng đa thức cho các nghiệm bội. Nếu $|z_\alpha| < 1$, các số hạng này suy giảm theo hàm mũ và bị chặn (sự tăng trưởng đa thức bị lấn át bởi $|z_\alpha|^n \to 0$). Nếu $|z_\alpha| = 1$ và là nghiệm đơn ($m_\alpha = 1$), thì $y_n = z_\alpha^n$ có $|y_n| = 1$, do đó vẫn bị chặn. Các hàm cơ sở này độc lập tuyến tính, tạo thành một tập cơ bản thông qua ma trận Vandermonde không kỳ dị cho các nghiệm phân biệt (hoặc các đạo hàm của nó cho nghiệm bội), vì vậy bất kỳ nghiệm nào $y_n = \sum c_\alpha v_{\alpha, n}$ đều thỏa mãn $|y_n| \le C$ với một hằng số $C$ nào đó không phụ thuộc vào $n$. Tính bị chặn cũng có thể được thiết lập bằng cách sử dụng định lý hình tròn Gershgorin trên ma trận đồng hành (\textit{companion matrix}) của hệ thức truy hồi, khẳng định các giá trị riêng nằm bên trong hoặc trên hình tròn đơn vị và không có khối Jordan nào lớn hơn kích thước $1 \times 1$ trên biên.
%\noindent \textbf{Mối ràng buộc chặt chẽ của sự hội tụ đối với sai số của các bước mồi khởi động}

%\noindent Khác với phương pháp một bước (chỉ cần xuất phát từ duy nhất $y_0 = y(t_0)$), công thức đa bước $k$-bước đòi hỏi phải được cung cấp sẵn một tập $k$ giá trị khởi động:
%\[
%    Y_0 = \{y_0, y_1, y_2, \dots, y_{k-1}\}.
%\]
%\noindent Do đó, sự hội tụ toàn cục của thuật toán không chỉ phụ thuộc vào cấu trúc nội tại của phương pháp đa bước mà còn chịu sự chi phối mang tính quyết định từ độ chính xác của quá trình chuẩn bị dữ liệu mồi này.

%\begin{itemize}
%    \item \textbf{Điều kiện tiệm cận ban đầu:} Để phương pháp đa bước có thể hội tụ, trước hết tập các giá trị xuất phát (thường được tạo ra từ một phương pháp Runge--Kutta) bắt buộc phải hội tụ về nghiệm đúng tại các thời điểm tương ứng khi bước lưới tiến dần đến không:
%    \begin{equation}
%        \lim_{h \to 0} \|Y(t_j) - y_j\| = 0, \qquad \forall j = 0, 1, \dots, k-1. \tag{3.11}
%    \end{equation}
%    \noindent Nếu có bất kỳ sự sai lệch nào ở bước mồi không thuyên giảm khi $h \to 0$, sai số đó sẽ bị khuếch đại theo chuỗi lặp đa bước và làm phương pháp mất tính hội tụ.

%    \item \textbf{Ràng buộc đồng bộ về bậc chính xác:} Giả sử phương pháp đa bước tuyến tính có bậc chính xác lý thuyết là $p \ge 1$ (tức sai số chặt cụt cục bộ tại mỗi bước thỏa mãn $d_{n+k} = \mathcal{O}(h^{p+1})$). Để toàn bộ quá trình tính toán đạt được đúng \textit{bậc hội tụ toàn cục $\mathcal{O}(h^p)$}, sai số ban đầu của các bước khởi động đòi hỏi phải bị chặn bởi một đại lượng cùng cấp:
%    \begin{equation}
%        \max_{0 \le j \le k-1} \|y(t_j) - y_j\| \le C_0 h^p, \tag{3.12}
%    \end{equation}
%    \noindent với $C_0 > 0$ là hằng số độc lập với $h$.
%    \item \textit{Hệ quả thực tiễn:} Nếu ta sử dụng một thuật toán mồi có bậc chính xác thấp hơn $p$ (chẳng hạn phương pháp Euler bậc $1$ mồi cho công thức Adams bậc $4$), sai số lớn từ giai đoạn khởi động ($\sim \mathcal{O}(h)$) sẽ ``kế thừa'' vào sơ đồ lặp và thống trị toàn bộ sai số toàn cục phía sau. Khi đó, phương pháp đa bước sẽ bị suy giảm bậc hội tụ nghiêm trọng, không thể phát huy được tiềm năng chính xác cao lý thuyết của nó.
%\end{itemize}

\subsubsection{Sự hội tụ}
%\begin{itemize}
%    \item \textbf{Định nghĩa toán học:} Nghiệm số $y_n$ tiệm cận chính xác về nghiệm giải tích $y(t)$ trên toàn miền khảo sát khi bước lưới $h \to 0$.
%    \item Mối ràng buộc chặt chẽ của sự hội tụ đối với sai số của các bước mồi khởi động.
%\end{itemize}

Mục tiêu cuối cùng của bất kỳ thuật toán giải tích số nào là đảm bảo nghiệm số thu được phải tiệm cận chính xác về nghiệm giải tích đích thực của phương trình vi phân khi ta tiến hành làm mịn lưới thời gian.

\vspace{0.5em}

\begin{baitapbox}{Sự hội tụ của phương pháp số}
    
    Phương pháp tuyến tính đa bước \eqref{eq:2.2} được gọi là \textbf{hội tụ} nếu, với mọi bài toán giá trị ban đầu (IVP) thỏa mãn các giả thiết của định lý Picard, ta có
    \begin{equation}
        \lim_{\substack{h \to 0 \\ n h = t - t_0}} Y_n = y(t_n). \tag{2.16} \label{eq:2.16}
    \end{equation}
    đúng với mọi $t \in [t_0, T_{max}]$ và với mọi nghiệm $\{y_n\}_{n=0}^N$ của phương trình sai phân \eqref{eq:2.2} với \textbf{các điều kiện khởi đầu nhất quán}, tức là với các điều kiện khởi đầu $y_s = \eta_s(h)$, $s = 0, 1, \dots, k - 1$, sao cho $\displaystyle\lim_{h \to 0} \eta_s(h) = y_0$, $s = 0, 1, \dots, k - 1$.\\
    
    % \noindent Nếu xét trên toàn bộ đoạn $[t_0, T]$, điều kiện hội tụ tương đương với việc chuẩn cực đại của sai số toàn cục tiến về không:
    % \[
        % \lim_{h \to 0} \max_{0 \le n \le N} \|Y(t_n) - y_n\| = 0.
    % \]
\end{baitapbox}

Ở đây, định nghĩa trên đòi hỏi điều kiện \eqref{eq:2.16} \textit{không những} phải đúng với các dãy $\{y_n\}_{n=0}^N$ được tạo ra từ \eqref{eq:2.2} bằng cách sử dụng các giá trị khởi đầu \textit{chính xác} $y_s = y(t_s)$, $s = 0, 1, \dots, k - 1$, mà còn phải đúng với mọi dãy $\{y_n\}_{n=0}^N$ có các giá trị khởi đầu $\eta_s(h)$ tiến tới giá trị đúng, $y_0$, khi $h \to 0$. Giả thiết này được đưa ra vì trong thực tế, các giá trị khởi đầu chính xác thường không có sẵn và phải được tính toán bằng phương pháp số.

\vspace{0.5em}

\begin{baitapbox}{Bậc hội tụ}
    Để đánh giá một phương pháp số hội tụ nhanh hay chậm về nghiệm chính xác của nó thì ta dựa trên \textit{bậc hội tụ} của phương pháp số đó. Phương pháp số hội tụ với bậc \textit{$p$} nếu tồn tại một hằng số $C > 0$ không phụ thuộc vào $h$ và $n$ sao cho:
   \[
    \max_{0 \le n\le N}|y(t_n) - Y_h(t_n)| \le Ch^p = \mathcal{O}(h^p)
   \]
   Bậc hội tụ càng cao thì phương pháp hội tụ càng nhanh.
\end{baitapbox}

\subsubsection{Định lý tương đương Dahlquist}
%\begin{itemize}
%    \item \textbf{Phát biểu định lý:} Đối với một phương pháp đa bước tuyến tính đặt chỉnh, điều kiện cần và đủ để phương pháp đạt sự \textbf{Hội tụ} là nó phải đồng thời thỏa mãn tính \textbf{Nhất quán} và tính \textbf{Ổn định}.
%    \item Khẳng định triết lý toán học số: $\text{Hội tụ} \Longleftrightarrow \text{Nhất quán} + \text{Ổn định}$.
%\end{itemize}

%Trong lý thuyết của các phương pháp giải tích số, sự nhất quán (Consistency), tính ổn định (Stability) và sự hội tụ (Convergence) là ba trụ cột nền tảng quyết định sự thành bại của một thuật toán. 

Để đánh giá một phương pháp giải tích số có nên được sử dụng hay không thì ta thường dựa vào ba yếu tố chính: \textit{Tính nhất quán}, \textit{tính ổn định} và  \textit{sự hội tụ} của phương pháp số về nghiệm chính xác. Mối liên hệ bản chất giữa ba đặc trưng này được thống nhất trọn vẹn thông qua một định lý kinh điển do nhà toán học Germund Dahlquist thiết lập vào năm 1956.

\vspace{0.5em}

Xét một bài toán giá trị đầu (IVP) được đặt chỉnh (well-posed) thỏa mãn điều kiện Lipschitz đảm bảo sự tồn tại và duy nhất nghiệm trên khoảng $[t_0, T]$. Khi ta áp dụng một phương pháp đa bước tuyến tính $k$-bước để xấp xỉ bài toán này:

\begin{baitapbox}{Định lý tương đương Dahlquist}
    Đối với một phương pháp đa bước tuyến tính với các hệ số hằng số đã được chuẩn hóa ($\alpha_k \neq 0, |\alpha_0| + |\beta_0| > 0$) và được xuất phát bằng một tập dữ liệu ban đầu tương thích, một phương pháp hội tụ khi và chỉ khi phương pháp đó \textit{nhất quán} và \textit{ổn định}.
    \[
        \textbf{Hội tụ} \iff \textbf{Nhất quán} + \textbf{Ổn định}
    \]
\end{baitapbox}

\vspace{0.5em}
%\noindent \textbf{2. Khẳng định triết lý toán học số}

%\noindent Định lý Dahlquist khẳng định một triết lý cốt lõi và mang tính chuẩn mực trong tính toán số:

%\noindent Triết lý này đem lại những hệ quả nhận thức và thực tiễn sâu sắc:
%\begin{itemize}
%   \item \textbf{Sự phân tách nhiệm vụ (Phân tích chi phí -- hiệu quả):} Định lý cho phép tách biệt việc kiểm tra tính hội tụ (vốn rất khó chứng minh trực tiếp theo định nghĩa do phụ thuộc vào nghiệm giải tích chưa biết) thành hai bài toán đại số đơn giản hơn nhiều:
%    \begin{enumerate}
%        \item Kiểm tra \textbf{tính Nhất quán}: Chỉ cần xác minh hai điều kiện tuyến tính trên hệ số đa thức đặc trưng là $\rho(1) = 0$ và $\rho'(1) = \sigma(1)$ (tương đương với phương pháp đạt bậc chính xác $p \ge 1$).
%        \item Kiểm tra \textbf{tính Ổn định}: Chỉ cần giải phương trình đa thức $\rho(\zeta) = 0$ và kiểm tra xem các nghiệm $\zeta_i$ có thỏa mãn điều kiện nghiệm Dahlquist (nằm trong hình tròn đơn vị và nghiệm trên biên là nghiệm đơn) hay không.
%    \end{enumerate}
%    \item 


% \textbf{Ý nghĩa trực quan về mặt cấu trúc:}
    % \begin{itemize}
        % \item \textit{Tính nhất quán} đóng vai trò điều khiển \textit{sai số cục bộ} tại mỗi bước lặp, bảo đảm phương pháp đi "đúng hướng" ở quy mô vi mô khi bước lưới $h \to 0$.
        % \item \textit{Tính ổn định} đóng vai trò như một cơ chế kiềm chế sự lan truyền nhiễu ở quy mô vĩ mô, bảo đảm rằng tổng các sai số cục bộ tích lũy qua $N \sim \mathcal{O}(h^{-1})$ bước lặp không bị khuếch đại bùng nổ.
        % \item Khi và chỉ khi phương pháp vừa "mô phỏng đúng cục bộ" (Nhất quán) vừa "kiểm soát được nhiễu toàn cục" (Ổn định), thì nghiệm xấp xỉ mới thực sự tiệm cận về nghiệm đúng chính xác (Hội tụ).
    % \end{itemize}
%\end{itemize}


%\subsubsection{Phác thảo chứng minh định lý hội tụ chính}
%\begin{itemize}
%    \item Chứng minh chiều thuận ($\Rightarrow$) bằng phản chứng.
%    \item Chứng minh chiều nghịch ($\Leftarrow$) dựa trên việc thiết lập phương trình sai số toàn cục và áp dụng bổ đề Gronwall dạng rời rạc để chặn trên sai số.
%\end{itemize}

%Để làm sáng tỏ bản chất toán học của Định lý tương đương Dahlquist, ta tiến hành phân tích và chứng minh chặt chẽ theo hai chiều suy luận tương đương.


\vspace{0.5em}
% \noindent \textbf{Chiều thuận (Điều kiện cần): Hội tụ $\implies$ Nhất quán $+$ 0-Ổn định}

% \noindent \textit{Ý tưởng chủ đạo:} Giả sử phương pháp đa bước là hội tụ cho mọi bài toán giá trị đầu thỏa mãn điều kiện Lipschitz. Để chứng minh tính nhất quán và 0-ổn định, ta thiết lập các ``bài toán thử''  đơn giản nhất có nghiệm giải tích tường minh. Nếu phương pháp đã hội tụ cho mọi bài toán thì nó bắt buộc phải hội tụ trên các bài toán thử này, từ đó áp đặt các ràng buộc đại số lên đa thức đặc trưng $\rho(\zeta)$ và $\sigma(\zeta)$.

% \textbf{Bước 1: Chứng minh tính nhất quán ($\rho(1) = 0$ và $\rho'(1) = \sigma(1)$)}
    
% \textit{Ràng buộc $\rho(1) = 0$:} Xét bài toán tầm thường $y'(t) = 0$ với điều kiện ban đầu $y(0) = 1$ trên $[0, T]$. Nghiệm chính xác hiển nhiên là $y(t) \equiv 1$. \\
% Khi áp dụng phương pháp đa bước với bước mồi chính xác $y_0 = y_1 = \dots = y_{k-1} = 1$, vế phải $f(t, y) = 0$, công thức trở thành:
% \[
%     \sum_{j=0}^{k} \alpha_j y_{n+j} = 0.
% \]
% Do phương pháp hội tụ về nghiệm hàm hằng $y(t) \equiv 1$ khi $h \to 0$, thay $y_{n+j} = 1$ vào phương trình sai phân, ta thu được:
% \[
%     \sum_{j=0}^{k} \alpha_j = 0 \iff \rho(1) = 0.
% \]
% \textit{Ràng buộc $\rho'(1) = \sigma(1)$:} Xét bài toán $y'(t) = 1$ với $y(0) = 0$. Nghiệm chính xác là $y(t) = t$. \\
%         Áp dụng phương pháp với nghiệm xấp xỉ hội tụ về nghiệm đúng $y_n = t_n = nh$ và $f(t_n, y_n) = 1$, ta có:
% \[
%     \sum_{j=0}^{k} \alpha_j (n+j)h = h \sum_{j=0}^{k} \beta_j \implies h n \sum_{j=0}^{k} \alpha_j + h \sum_{j=0}^{k} j \alpha_j = h \sum_{j=0}^{k} \beta_j.
% \]
% Do $\sum_{j=0}^k \alpha_j = \rho(1) = 0$, chia hai vế cho $h \neq 0$, ta thu được $\sum_{j=0}^k j \alpha_j = \sum_{j=0}^k \beta_j$, hay nói cách khác: $\rho'(1) = \sigma(1)$. Vậy phương pháp chắc chắn phải nhất quán.\\
% \textbf{Bước 2: Chứng minh tính 0-Ổn định (Thỏa mãn Điều kiện nghiệm Dahlquist)}\\
% Tiếp tục xét bài toán thử $y'(t) = 0, y(0) = 0$ với nghiệm đúng $y(t) \equiv 0$. Phương trình sai phân tương ứng là $\sum_{j=0}^k \alpha_j y_{n+j} = 0$.\\
% Giả sử phản chứng rằng phương pháp \textit{không} 0-ổn định. Theo Định lý Dahlquist về phương trình sai phân, đa thức $\rho(\zeta) = 0$ phải có ít nhất một nghiệm vi phạm điều kiện nghiệm, tức là:
%         \begin{enumerate}
%             \item Hoặc tồn tại nghiệm $\zeta_s$ với $|\zeta_s| > 1$. Khi đó nghiệm sai phân có dạng $y_n = h \zeta_s^n$. Dù sai số ban đầu $y_0, \dots, y_{k-1} \sim \mathcal{O}(h) \to 0$, nhưng khi $n = \frac{t}{h} \to \infty$, giá trị $|y_n| \to \infty$, mâu thuẫn với sự hội tụ về $0$.
%             \item Hoặc tồn tại nghiệm bội $m \ge 2$ trên đường tròn đơn vị ($|\zeta_m| = 1$). Khi đó phương trình sai phân có nghiệm riêng dạng $y_n = h \cdot n^{m-1} \zeta_m^n$. Tại thời điểm $t_n = nh = t > 0$ cố định, ta có $y_n = t^{m-1} h^{2-m} \zeta_m^n$. Vì $m \ge 2$, đại lượng này không thể hội tụ về $0$ khi $h \to 0$.
%         \end{enumerate}
% Do đó, giả thiết phản chứng là sai. Phương pháp bắt buộc phải 0-ổn định.
   
% \vspace{0.8em}

% \noindent \textbf{Chiều đảo (Điều kiện đủ): Nhất quán $+$ 0-Ổn định $\implies$ Hội tụ}

% \noindent \textit{Ý tưởng chủ đạo:} Ta chuyển phương trình sai phân đa bước vô hướng cấp $k$ về một hệ phương trình sai phân vectơ cấp $1$. Sau đó, sử dụng tính 0-ổn định để bị chặn chuẩn của ma trận lũy thừa, kết hợp với tính Nhất quán (sai số cục bộ nhỏ) và áp dụng \textbf{Bổ đề Gronwall rời rạc} để chứng minh sai số toàn cục $e_n \to 0$.

% \textbf{Bước 1: Thiết lập phương trình sai số toàn cục} \\
% Gọi $e_n = y(t_n) - y_n$ là sai số toàn cục. Trừ phương trình nghiệm số khỏi phương trình nghiệm đúng (đã chứa sai số chặt cụt cục bộ $h \alpha_k d_{n+k}$), ta có phương trình lan truyền sai số:
% \begin{equation}
%     \sum_{j=0}^{k} \alpha_j e_{n+j} = h \sum_{j=0}^{k} \beta_j \left[ f(t_{n+j}, y(t_{n+j})) - f(t_{n+j}, y_{n+j}) \right] + h \alpha_k d_{n+k}. \tag{3.14}
% \end{equation}
% Áp dụng điều kiện Lipschitz $\|f(t, u) - f(t, v)\| \le L\|u - v\|$, đặt $g_{n+j} = f(t_{n+j}, y(t_{n+j})) - f(t_{n+j}, y_{n+j})$, ta có $\|g_{n+j}\| \le L \|e_{n+j}\|$. Chuẩn hóa $\alpha_k = 1$.

% \textbf{Bước 2: Chuyển về dạng hệ sai phân bậc nhất (Đồng hành Ma trận)} \\
% Định nghĩa vectơ sai số trạng thái $k$ chiều $E_n = [e_{n+k-1}, e_{n+k-2}, \dots, e_n]^T$. Khi đó, phương trình vô hướng (3.14) được viết gọn dưới dạng ma trận:
% \begin{equation}
%     E_{n+1} = A E_n + h \Phi_n + h D_n, \tag{3.15}
% \end{equation}
% trong đó $A$ là ma trận đồng hành (Companion matrix) cấu tạo từ các hệ số $\alpha_j$:
% \[
%     A = \begin{bmatrix}
%     -\alpha_{k-1} & -\alpha_{k-2} & \dots & -\alpha_1 & -\alpha_0 \\
%     1 & 0 & \dots & 0 & 0 \\
%     0 & 1 & \dots & 0 & 0 \\
%     \vdots & \vdots & \ddots & \vdots & \vdots \\
%     0 & 0 & \dots & 1 & 0
%     \end{bmatrix}.
% \]

% \textbf{Bước 3: Khai thác tính 0-Ổn định} \\
% Các giá trị riêng của ma trận $A$ chính là các nghiệm của đa thức đặc trưng $\rho(\zeta) = 0$. Vì phương pháp là \textbf{0-ổn định}, mọi giá trị riêng đều thỏa mãn $|\lambda_i| \le 1$ và các khối Jordan ứng với trị riêng trên đường tròn đơn vị đều có cấp $1$. \\
% Theo định lý đại số tuyến tính, tồn tại một chuẩn ma trận $\|\cdot\|$ sao cho $\|A^n\| \le K$ với mọi $n \ge 0$, trong đó $K > 0$ là hằng số độc lập với $n$ và $h$.\\

% \textbf{Bước 4: Áp dụng Bổ đề Gronwall rời rạc và Kết luận} \\
% Lấy chuẩn hai vế phương trình lặp ma trận (3.15) và sử dụng tính bị chặn $\|A^n\| \le K$, kết hợp với điều kiện Lipschitz, ta thiết lập được bất phương trình sai phân:
% \[
%     \|E_{n+1}\| \le \|E_n\| + h K L \beta^* \max_{0 \le j \le k} \|E_{n+j-1}\| + K h \|d_{\max}\|,
% \]
% với $\beta^* = \sum |\beta_j|$. Theo \textbf{Bổ đề Gronwall rời rạc}, sai số toàn cục tại nút $t_N = T$ bị chặn bởi:
%     \begin{equation}
%         \|e_N\| \le \|E_N\| \le K^* \left( \|E_0\| + \max_{0 \le n \le N} \|d_n\| \right) e^{L^* (T - t_0)}, \tag{3.16}
%     \end{equation}
%     trong đó $K^*, L^*$ là các hằng số không phụ thuộc vào $h$. 

% \textbf{Kết luận:} Do phương pháp \textbf{nhất quán}, sai số chặt cụt cục bộ $\max \|d_n\| \to 0$ khi $h \to 0$. Đồng thời, do bước khởi động chính xác nên sai số mồi $\|E_0\| \to 0$. Từ bất đẳng thức (3.16), ta áp dụng định lý kẹp và suy ra $\lim_{h \to 0} \|e_N\| = 0$. Phương pháp hội tụ hoàn toàn. $\blacksquare$

\subsubsection{Chứng minh tương đương Dahlquist}
\textbf{Chiều thuận} (Hội tụ $\Longrightarrow$ Nhất quán + $0$-ổn định)

Trước hết ta đi chứng minh phương pháp đa bước tuyến tính hội thì có tính $0$-ổn đinh. Ta xét bài toán giá trị ban đầu $y' = 0$, $y(0) = 0$, trên khoảng $[0, T_{max}]$, $T_{max} > 0$, bài toán này có là $y(t) \equiv 0$. Thay vào công thức phương pháp đa bước, ta được phương trình sai phân:
\begin{equation}
    \alpha_k y_{n+k} + \alpha_{k-1} y_{n+k-1} + \dots + \alpha_0 y_n = 0 . \quad \tag{2.16} \label{eq:2.16}
\end{equation}

Vì phương pháp được giả định là hội tụ, với bất kỳ $t > 0$ nào, ta có:
\begin{equation}
    \lim_{\substack{h \to 0 \\ nh=t}} Y_n = 0 , \quad \tag{2.17} \label{eq:2.17}
\end{equation}
đối với mọi nghiệm của \eqref{eq:2.16} thỏa mãn $y_s = \eta_s(h)$, $s = 0, \dots, k - 1$, trong đó:
\begin{equation}
    \lim_{h \to 0} \eta_s(h) = 0 , \quad s = 0, 1, \dots k - 1 . \quad \tag{2.18} \label{eq:2.18}
\end{equation}
    

Gọi $z = r e^{i\phi}$, là một nghiệm của đa thức đặc trưng thứ nhất $\rho(z)$; $r \ge 0$, $0 \le \phi < 2\pi$. Có thể dễ dàng kiểm chứng rằng các số:
\[
    y_n = h r^n \cos n\phi
\]
xác định một nghiệm cho \eqref{eq:2.16} thỏa mãn \eqref{eq:2.18}). Nếu $\phi \neq 0$ và $\phi \neq \pi$, thì:
\[
    \frac{y_n^2 - y_{n+1}y_{n-1}}{\sin^2 \phi} = h^2 r^{2n} .
\]

Vì vế trái của đồng nhất thức này hội tụ về $0$ khi $h \to 0, n \to \infty, nh = t$, điều tương tự cũng phải đúng cho vế phải; do đó:
\[
    \lim_{n \to \infty} \left(\frac{x}{n}\right)^2 r^{2n} = 0 \Leftrightarrow r \le 1
\]

Vậy ta đã chứng minh được rằng mọi nghiệm của đa thức đặc trưng thứ nhất của \eqref{eq:2.2} đều nằm trong hình tròn đơn vị đóng.

Tiếp theo ta chứng minh rằng bất kỳ nghiệm nào của đa thức đặc trưng thứ nhất của \eqref{eq:2.2} nằm trên đường tròn đơn vị đều phải là \textit{nghiệm đơn}. Giả sử ngược lại, $z = r e^{i\phi}$, là một nghiệm bội của $\rho(z)$, với $|z| = 1$ (và do đó $r = 1$) và $0 \le \phi < 2\pi$. Có thể dễ dàng kiểm tra rằng các số:
\begin{equation}
     y_n = h^{1/2} n r^n \cos(n\phi) \quad \tag{2.19}\label{eq:2.19}
\end{equation}
xác định một nghiệm cho \eqref{eq:2.16} thỏa mãn \eqref{eq:2.18}, vì:
\[
    |\eta_s(h)| = |y_s| \le h^{1/2} s \le h^{1/2}(k - 1), \quad s = 0, \dots k - 1 .
\]

Nếu $\phi = 0$ hoặc $\phi = \pi$, từ \eqref{eq:2.19} với $h = x/n$ suy ra rằng:
\begin{equation}
    |y_n| = x^{1/2} n^{1/2} r^n . \quad \tag{2.20} \label{eq:2.20}
\end{equation}

Vì theo giả thiết, $|z| = 1$ (và do đó $r = 1$), ta suy ra từ \eqref{eq:2.20} rằng $\displaystyle \lim_{n \to \infty} |Y_n| = \infty$, điều này mâu thuẫn với \eqref{eq:2.17}.

Mặt khác, nếu $\phi \neq 0$ và $\phi \neq \pi$, thì:
\begin{equation}
    \frac{z_n^2 - z_{n+1}z_{n-1}}{\sin^2 \phi} = r^{2n} , \quad \tag{2.21} \label{eq:2.21}
\end{equation}
trong đó $z_n = n^{-1}h^{-1/2}y_n = h^{1/2}x^{-1}y_n$. Vì theo \eqref{eq:2.17}, $\lim_{n \to \infty} z_n = 0$, suy ra vế trái của \eqref{eq:2.21} hội tụ về $0$ khi $n \to \infty$. Tuy nhiên, vế trái của \eqref{eq:2.17} không thể hội tụ về $0$ khi $n \to \infty$, vì $|z| = 1$ (và do đó $r = 1$). Do đó, một lần nữa, ta lại đi đến một mâu thuẫn.

Tóm lại, ta đã chứng minh được rằng tất cả các nghiệm của đa thức đặc trưng thứ nhất $\rho(z)$ của phương pháp đa bước tuyến tính \eqref{eq:2.2} đều nằm trong hình tròn đơn vị $|z| \le 1$, và những nghiệm thuộc đường tròn đơn vị $|z| = 1$ đều là nghiệm đơn. Dựa vào định lý về tính $0$-ổn định thì phương pháp đa bước tuyến tính này $0$-ổn định.\\


Tiếp đến ta đi chứng minh Hội tụ $\Longrightarrow$ nhất quán. Ta xét bài toán (IVP):
\[
y'=0, \quad y(0)=1 
\]

Bài toán có nghiệm $y(t) = 1$ trên đoạn $[0;T_{max}] \ (T_{max} >0)$. Thay vào công thức của phương pháp đa bước tuyến tính, ta được:
\[
\sum_{j=0}^{k} \alpha_j Y_{n+j} = h \sum_{j=0}^{k} \beta_j f(t_{n+j}, y(t_{n+j})) = h \sum_{j=0}^{k}\beta_j0=0
\]
\begin{equation}
    \Leftrightarrow \alpha_kY_{n+k} + \alpha_{k-1}Y_{n+k-1} + \dots + \alpha_{0}Y_{n} = 0 \tag{2.22} \label{eq:2.22}
\end{equation}

Không mất tính tổng quát, ta chọn các giá trị mồi $Y_k = 1 \ \forall k = \overline{0,..,k-1}$. Kết hợp với tính hội tụ thì cho ta:
\[
\lim_{h \to 0 \\ nh = t} Y_n = y(t) = 1
\]
\[
\Leftrightarrow \lim_{n \to \infty} Y_n = 1
\]

Ta lấy giới hạn hai vế \eqref{eq:2.22} cho $n \to \infty$, thì được:
\[
\alpha_{k} + \alpha_{k-1} + \dots + \alpha_{0} = 0 
\]

Để ý vế phải là $\rho(1)$, vậy ta có:
\begin{equation}
    \rho(1) = 0 \tag{2.23} \label{eq:2.23}
\end{equation}

Xét bài toán (IVP):
\[
y' = 1 \qquad y(0) = 0 
\]

Bài toán có nghiệm $y(t) = x$ trên đoạn $[0;T_{max}] \ (T_{max} >0)$.Tương tự trên, ta thay vào công thức phương pháp đa bước thì được:
\[
\alpha_{k}Y_{n+k} + \alpha_{k-1}Y_{n+k-1} + \dots + \alpha_{0}Y_{n} = h\left(\beta_{k}+\beta_{k-1}+ \dots + \beta_{0}\right)
\]

Giả sử phương trình trên có nghiệm $Y_s, \ s=\overline{0,\dots,k-1}$. Với một phương pháp hội tụ thì phải đồng thời thỏa mãn:
\[
\lim_{h \to 0} Y_s = 0, \quad s=\overline{0,\dots,k-1}, \qquad \lim_{h \to 0\\ nh = t}Y_n = y(t) = x
\]

Với chứng minh trên, ta có một phương pháp hội tụ thì có tính $0$-ổn định. Điều này tương đương với mọi nghiệm $\zeta_j$ của đa thức đặc trưng $\rho(\zeta)$ đều thỏa mãn điều kiện nghiệm. Với $\zeta = 1$ là nghiệm của đa thức đặc trưng trên thì $\zeta = 1$ là nghiệm đơn, tức là:
\[
\rho'(1) = k\alpha_{k} + ...+2\alpha_{2} + \alpha_{1} \ne 0
\]

Tương tự trên ta gán các giá trị mồi bằng dãy $\{Y_n\}^N_{n=1}$ định nghĩa bởi $Y_n =Knh$ với:
\[
K = \frac{\rho(1)}{\rho'(1)} = \frac{\beta_{k}+\dots+\beta_{1}+\beta_{0}}{k\alpha_{k}+\dots+2\alpha_{2}+\alpha_{1}}
\]

Dãy trên đồng thời thỏa mãn điều kiện (2.). Do đó ta có:
\[
t = y(t) = \lim_{h \to 0 \\ nh = t }Y_n = \lim_{h \to 0 \\ nh = t}Knh = Kt
\]

Từ đây ta có $K=1$ tương đương với:
\begin{equation}
    \rho'(1) = \sigma(1) \tag{2.24} \label{eq:2.24}
\end{equation}

Từ \eqref{eq:2.23} và \eqref{eq:2.24} ta có phương pháp đa bước hội tụ thì có tính nhất quán.\\

\noindent Trước khi đi vào chứng minh chiều đảo ta công nhận hai bổ đề sau:

\begin{baitapbox}{Bổ đề 1}
    Giả sử rằng tất cả các nghiệm của đa thức $\rho(z) = \alpha_k z^k + \cdots + \alpha_1 z + \alpha_0$ nằm trong hình tròn đơn vị đóng $|z| \le 1$ và những nghiệm nằm trên đường tròn đơn vị $|z| = 1$ là các nghiệm đơn. Giả sử thêm rằng các số $\gamma_l$, với $l = 0, 1, 2, \dots$, được xác định bởi
\begin{equation*}
    \frac{1}{\alpha_k + \cdots + \alpha_1 z^{k-1} + \alpha_0 z^k} = \gamma_0 + \gamma_1 z + \gamma_2 z^2 + \cdots .
\end{equation*}
Khi đó, $\Gamma \equiv \sup_{l \ge 0} |\gamma_l| < \infty$.
\end{baitapbox}
\noindent\textbf{Chứng minh:}\\

Trước hết ta định nghĩa: 
\begin{equation*}
    \hat{\rho}(z) = z^k \rho\left(\frac{1}{z}\right)
\end{equation*} 
và lưu ý rằng, nhờ các giả thiết về nghiệm của $\rho(z)$, hàm số ${1}/\hat{\rho}(z)$ là chỉnh hình trong hình tròn đơn vị mở $|z| < 1$. Vì các nghiệm $z_1, z_2, \dots, z_m$ của $\rho(z)$ trên $|z| = 1$ là nghiệm đơn, điều tương tự cũng đúng với các cực điểm của $\frac{1}{\hat{\rho}(z)}$, và tồn tại các hằng số $A_1, \dots, A_m$ sao cho hàm số
\begin{equation}
    f(z) = \frac{1}{\hat{\rho}(z)} - \frac{A_1}{z - \frac{1}{z_1}} - \cdots - \frac{A_m}{z - \frac{1}{z_m}} \tag{2.25} \label{eq:2.25}
\end{equation}
chỉnh hình với $|z| < 1$ và $|f(z)| \le M$ với mọi $|z| \le 1$. Do đó, theo ước lượng Cauchy\footnote{\textbf{Định lý (Ước lượng Cauchy)} Nếu $f$ là một hàm chỉnh hình trong hình tròn mở $D(a, R)$ tâm $a$ bán kính $R$, và $|f(z)| \le M$ với mọi $z \in D(a, R)$, thì $|f^{(n)}(a)| \le M(n! / R^n)$, $n = 0, 1, 2, \dots$. [Xem chứng minh của kết quả này, ví dụ, trong Walter Rudin: \textit{Real and Complex Analysis}. 3rd edition. McGraw-Hill, New York, 1986.]}, các hệ số của khai triển Taylor của $f$ tại $z = 0$ cũng tạo thành một dãy bị chặn. Vì
\begin{equation*}
    -\frac{A_i}{z - \frac{1}{z_i}} = A_i \sum_{l=0}^\infty z_i^l z^l, \quad i = 1, \dots, m,
\end{equation*}
và $|z_i| \le 1$, từ \eqref{eq:2.25} ta suy ra rằng các hệ số trong khai triển chuỗi Taylor của $1/\hat{\rho}(z)$ tạo thành một dãy bị chặn, điều này hoàn tất chứng minh.

Bây giờ ta sẽ áp dụng Bổ đề 1 để ước lượng nghiệm của phương trình sai phân tuyến tính
\begin{equation}
    \alpha_k e_{m+k} + \alpha_{k-1} e_{m+k-1} + \cdots + \alpha_0 e_0 = h(\beta_{k,m} e_{m+k} + \beta_{k-1,m} e_{m+k-1} + \cdots + \beta_{0,m} e_m) + \lambda_m . \tag{2.26} \label{eq:2.26}
\end{equation}



\begin{baitapbox}{Bổ đề 2}
    Giả sử rằng tất cả các nghiệm của đa thức $\rho(z) = \alpha_k z^k + \cdots + \alpha_1 z + \alpha_0$ nằm trong hình tròn đơn vị đóng $|z| \le 1$ và những nghiệm nằm trên đường tròn đơn vị $|z| = 1$ là các nghiệm đơn. Gọi $B^*$ và $\Lambda$ là các hằng số không âm và $\beta$ là một hằng số dương sao cho
\begin{equation*}
    |\beta_{k,n}| + |\beta_{k-1,n}| + \cdots + |\beta_{0,n}| \le B^*,
\end{equation*}
\begin{equation*}
    |\beta_{k,n}| \le \beta, \quad |\lambda_n| \le \Lambda, \quad n = 0, 1, \dots, N,
\end{equation*}
và giả sử $0 \le h < |\alpha_k|\beta^{-1}$. Khi đó, mọi nghiệm của \eqref{eq:2.26} thỏa mãn
\begin{equation*}
    |e_s| \le E, \quad s = 0, 1, \dots, k - 1,
\end{equation*}
sẽ thỏa mãn
\begin{equation*}
    |e_n| \le K^* \exp(nhL^*), \quad n = 0, 1, \dots, N,
\end{equation*}
trong đó
\begin{equation*}
    L^* = \Gamma^* B^*, \quad K^* = \Gamma^*(N\Lambda + AEk), \quad \Gamma^* = \Gamma / (1 - h|\alpha_k|^{-1}\beta),
\end{equation*}
$\Gamma$ được xác định như trong Bổ đề 1, và
\begin{equation*}
    A = |\alpha_k| + |\alpha_{k-1}| + \cdots + |\alpha_0|.
\end{equation*}
\end{baitapbox}

Với một $k$ cố định, ta xét các số $\gamma_l$, $l = 0, 1, \dots, n - k$, được định nghĩa trong Bổ đề 1. Sau khi nhân cả hai vế của phương trình \eqref{eq:2.26} ứng với $m = n - k - l$ cho $\gamma_l$, $l = 0, 1, \dots, n - k$ và cộng các phương trình thu được, ký hiệu tổng là $S_n$, bằng cách biến đổi vế trái của tổng ta tìm được:
\begin{align*}
S_n &= (\alpha_k e_n + \alpha_{k-1} e_{n-1} + \cdots + \alpha_0 e_{n-k})\gamma_0 \\
&\quad + (\alpha_k e_{n-1} + \alpha_{k-1} e_{n-2} + \cdots + \alpha_0 e_{n-k-1})\gamma_1 + \cdots \\
&\quad + (\alpha_k e_k + \alpha_{k-1} e_{k-1} + \cdots + \alpha_0 e_0)\gamma_{n-k} \\
&= \alpha_k \gamma_0 e_n + (\alpha_k \gamma_1 + \alpha_{k-1} \gamma_0)e_{n-1} + \cdots \\
&\quad + (\alpha_k \gamma_{n-k} + \alpha_{k-1} \gamma_{n-k-1} + \cdots + \alpha_0 \gamma_{n-2k})e_k \\
&\quad + (\alpha_{k-1} \gamma_{n-k} + \cdots + \alpha_0 \gamma_{n-2k+1})e_{k-1} + \cdots \\
&\quad + \alpha_0 \gamma_{n-k} e_0 .
\end{align*}

Định nghĩa $\gamma_l = 0$ với $l < 0$ và chú ý rằng
\begin{equation}
    \alpha_k \gamma_l + \alpha_{k-1} \gamma_{l-1} + \cdots + \alpha_0 \gamma_{l-k} = 
    \begin{cases} 
    1 & \text{với } l = 0 \\ 
    0 & \text{với } l > 0 
    \end{cases} \tag{2.27} \label{eq:2.27}
\end{equation}
ta có
\begin{equation*}
    S_n = e_n + (\alpha_{k-1} \gamma_{n-k} + \cdots + \alpha_0 \gamma_{n-2k+1})e_{k-1} + \cdots + \alpha_0 \gamma_{n-k} e_0 .
\end{equation*}

Bằng cách biến đổi tương tự vế phải trong tổng, ta thu được
\begin{align*}
&e_n + (\alpha_{k-1} \gamma_{n-k} + \cdots + \alpha_0 \gamma_{n-2k+1})e_{k-1} + \cdots + \alpha_0 \gamma_{n-k} e_0 \\
&\quad = h \big[ \beta_{k, n-k} \gamma_0 e_n + (\beta_{k-1, n-k} \gamma_0 + \beta_{k, n-k-1} \gamma_1)e_{n-1} + \cdots \\
&\quad\quad + (\beta_{0, n-k} \gamma_0 + \cdots + \beta_{k, n-2k} \gamma_k)e_{n-k} + \cdots + \beta_{0,0} \gamma_{n-k} e_0 \big] \\
&\quad\quad + (\lambda_{n-k} \gamma_0 + \lambda_{n-k-1} \gamma_1 + \cdots + \lambda_0 \gamma_{n-k}) .
\end{align*}

Do đó, theo các giả thiết và chú ý rằng theo \eqref{eq:2.27} ta có $\gamma_0 = \alpha_k^{-1}$,
\begin{equation*}
    |e_n| \le h\beta|\alpha_k^{-1}| |e_n| + h\Gamma B^* \sum_{m=0}^{n-1} |e_m| + N\Gamma\Lambda + A\Gamma E k .
\end{equation*}

\begin{equation*}
  \Rightarrow  (1 - h\beta|\alpha_k^{-1}|)|e_n| \le h\Gamma B^* \sum_{m=0}^{n-1} |e_m| + N\Gamma\Lambda + A\Gamma E k .
\end{equation*}

Nhắc lại định nghĩa của $L^*$ và $K^*$, ta có thể viết lại bất đẳng thức cuối cùng như sau:
\begin{equation}
    |e_n| \le K^* + h L^* \sum_{m=0}^{n-1} |e_m|, \quad n = 0, 1, \dots, N . \tag{2.28} \label{eq:2.28}
\end{equation}

Ước lượng cuối cùng được suy ra từ \eqref{eq:2.28} bằng quy nạp. Đầu tiên, ta chú ý rằng nhờ \eqref{eq:2.27}, $A\Gamma \ge 1$. Do đó, $K^* \ge \Gamma A E k \ge E k \ge E$. Bây giờ, cho $n = 1$ trong \eqref{eq:2.28},
\begin{equation*}
    |e_1| \le K^* + h L^* |e_0| \le K^* + h L^* E \le K^*(1 + h L^*) .
\end{equation*}

Lặp lại lập luận này, ta nhận thấy rằng
\begin{equation*}
    |e_m| \le K^*(1 + h L^*)^m, \quad m = 0, \dots, k - 1 .
\end{equation*}

Bây giờ giả sử bất đẳng thức này đã được chứng minh là đúng cho $m = 0, 1, \dots, n - 1$, với $n \ge k$; ta sẽ chứng minh rằng nó cũng đúng cho $m = n$, điều này sẽ hoàn tất phép chứng minh quy nạp. Thật vậy, từ \eqref{eq:2.28} ta có
\begin{equation}
    |e_n| \le K^* + h L^* K^* \frac{(1 + h L^*)^n - 1}{h L^*} = K^*(1 + h L^*)^n . \tag{2.29} \label{eq:2.29}
\end{equation}

Hơn nữa, vì $1 + h L^* \le e^{h L^*}$, từ \eqref{eq:2.29}\ ta có
\begin{equation}
    |e_n| \le K^* e^{h L^* n}, \quad n = 0, 1, \dots, N . \tag{2.30} \label{eq:2.30}
\end{equation}

Điều này hoàn tất chứng minh bổ đề. Ta lưu ý rằng phép suy ra \eqref{eq:2.28} $\Rightarrow$ \eqref{eq:2.30} thường được gọi là \textbf{Bổ đề Gronwall rời rạc (Discrete Gronwall Lemma)}. Từ đây ta tiến hành chứng minh chiều đảo:\\

\noindent\textbf{Chiều đảo} (Hội tụ $\Longleftarrow$ Nhất quán + $0$-ổn định)\\

Ta định nghĩa
\begin{equation*}
    \delta = \delta(h) = \max_{0 \le s \le k-1} |\eta_s(h) - y(a + sh)| ;
\end{equation*}
do các điều kiện khởi tạo $y_s = \eta_s(h)$, $s = 0, \dots, k$, được giả thiết là nhất quán, ta có $\displaystyle\lim_{h \to 0} \delta(h) = 0$. Ta cần chứng minh rằng
\begin{equation*}
    \lim_{\substack{n \to \infty \\ nh = t - t_0}} y_n = y(t)
\end{equation*}
với mọi $t$ thuộc khoảng $[t_0, T_{max}]$. Ta bắt đầu bài chứng minh bằng cách ước lượng sai số cắt cụt của \eqref{eq:2.2}:
\begin{equation}
    T_n = \frac{1}{h\sigma(1)} \left[ \sum_{j=0}^k \alpha_j y(t_{n+j}) - h \beta_j y'(t_{n+j}) \right] \qquad \tag{2.31} \label{eq:2.31}
\end{equation}

Vì $y' \in C[t_0, T_{max}]$, ta hoàn toàn có thể định nghĩa hàm số sau với $\epsilon \ge 0$:
\begin{equation*}
    \chi(\epsilon) = \max_{\substack{|t^* - t| \le \epsilon \\ t, t^* \in [t_0, T_{max}]}} |y'(t^*) - y'(t)| .
\end{equation*}

Với $s = 0, 1, \dots, k$, ta có thể viết
\begin{equation*}
    y'(t_{n+s}) = y'(t_n) + \theta_s \chi(sh) ,
\end{equation*}
trong đó $|\theta_s| \le 1$. Hơn nữa, theo Định lý Giá trị Trung bình (Mean Value Theorem), tồn tại $\xi_s \in (t_n, t_{n+s})$ sao cho
\begin{equation*}
    y(t_{n+s}) = y(t_n) + sh y'(\xi_s) .
\end{equation*}

Do đó,
\begin{equation*}
    y(t_{m+s}) = y(t_m) + sh \left[ y'(t_m) + \theta'_s \chi(sh) \right] ,
\end{equation*}
trong đó $|\theta'_s| \le 1$.

Bây giờ ta có thể viết
\begin{align*}
|\sigma(1)T_n| &\le \Big| h^{-1}(\alpha_1 + \alpha_2 + \cdots + \alpha_k)y(x_n) + (\alpha_1 + 2\alpha_2 + \cdots + k\alpha_k)y'(x_n) \\
&\quad - (\beta_0 + \beta_1 + \cdots + \beta_k)y'(x_n) \Big| \\
&\quad + (|\alpha_1| + 2|\alpha_2| + \cdots + k|\alpha_k|)\chi(kh) + (|\beta_0| + |\beta_1| + \cdots + |\beta_k|)\chi(kh) .
\end{align*}

Vì phương pháp đã được giả thiết là nhất quán, số hạng thứ nhất, thứ hai và thứ ba ở vế phải triệt tiêu lẫn nhau, cho ta
\begin{equation*}
    |\sigma(1)T_n| \le (|\alpha_1| + 2|\alpha_2| + \cdots + k|\alpha_k|)\chi(kh) + (|\beta_0| + |\beta_1| + \cdots + |\beta_k|)\chi(kh) .
\end{equation*}

Do đó,
\begin{equation}
    |\sigma(1)T_n| \le K \chi(kh) \qquad \tag{2.32} \label{eq:2.32}
\end{equation}
trong đó
\begin{equation*}
    K = |\alpha_1| + 2|\alpha_2| + \cdots + k|\alpha_k| + |\beta_0| + |\beta_1| + \cdots + |\beta_k| .
\end{equation*}

So sánh \eqref{eq:2.2} với \eqref{eq:2.31}, ta kết luận rằng sai số toàn cục $e_m = y(x_m) - Y_m$ thỏa mãn
\begin{equation*}
    \alpha_k e_{m+k} + \cdots + \alpha_0 e_0 = h(\beta_k g_{m+k} e_{m+k} + \cdots + \beta_0 g_m e_m) + \sigma(1) T_n h ,
\end{equation*}
trong đó
\begin{equation*}
    g_m = 
    \begin{cases}
    [f(x_m, y(x_m)) - f(x_m, y_m)] / e_m , & e_m \ne 0 \\
    0 , & e_m = 0 .
    \end{cases}
\end{equation*}

Nhờ có \eqref{eq:2.32}, ta suy ra
\begin{equation*}
    \alpha_k e_{m+k} + \cdots + \alpha_0 e_0 = h(\beta_k g_{m+k} e_{m+k} + \cdots + \beta_0 g_m e_m) + \theta K \chi(kh) h .
\end{equation*}

Vì $f$ được giả thiết là thỏa mãn điều kiện Lipschitz (3), ta có $|g_m| \le L$, $m = 0, 1, \dots$. Áp dụng Bổ đề 2 với $E = \delta(h)$, $\Lambda = K \chi(kh) h$, $N = (T_{max} - t_0)/h$, $B^* = BL$, trong đó $B = |\beta_0| + |\beta_1| + \cdots + |\beta_k|$, ta tìm được
\begin{equation}
    |e_n| \le \Gamma^* \left[ A k \delta(h) + (T_{max} - t_0) K \chi(kh) \right] \exp \left[ (t_n - t_0) L \Gamma^* B \right] , \tag{2.33} \label{eq:2.33}
\end{equation}
trong đó
\begin{equation*}
    A = |\alpha_0| + |\alpha_1| + \cdots + |\alpha_k| , \quad \Gamma^* = \frac{\Gamma}{1 - h|\alpha_k|^{-1}\beta_k| L} .
\end{equation*}

Bây giờ, do $y'$ là một hàm liên tục trên đoạn đóng $[t_0, T_{max}]$; do đó nó liên tục đều trên $[t_0, T_{max}]$. Vì vậy, $\chi(kh) \to 0$ khi $h \to 0$; đồng thời, nhờ giả thiết về tính nhất quán của các giá trị khởi tạo, $\delta(h) \to 0$ khi $h \to 0$. Chuyển sang giới hạn $h \to 0$ trong \eqref{eq:2.33}, ta suy ra rằng
\begin{equation*}
    \lim_{\substack{n \to \infty \\ t - t_0 = nh}} |e_n| = 0 ;
\end{equation*}
tương đương với,
\begin{equation*}
    \lim_{\substack{n \to \infty \\ t - t_0 = nh}} |y(t) - Y_n| = 0
\end{equation*}
vì vậy phương pháp là hội tụ.

